% Chapter 6 — Results and Discussion (SHAP, error analysis, limitations).
\section{Results and Discussion}

\subsection{Feature importance}

SHAP global importance (Figure~\ref{fig:shap}) is dominated by context overlap and answer length: HaluEval hallucinated answers tend to share vocabulary with the context while being systematically different in length from correct answers (median 10 vs.\ 2 words), an artifact of the benchmark's synthetic generation. NLI neutrality is the fifth most important signal. These results align with the ablation study and warn that benchmark-specific length patterns may not transfer.

\begin{figure}[t]
\centering
\includegraphics[width=0.62\textwidth]{fig_shap_importance.pdf}
\caption{Mean $|$SHAP$|$ feature importance (top 10).}
\label{fig:shap}
\end{figure}

\begin{figure}[t]
\centering
\includegraphics[width=0.85\textwidth]{fig_shap_summary.pdf}
\caption{SHAP beeswarm summary over the test set.}
\label{fig:shap2}
\end{figure}

\subsection{Error analysis}

The test set contains only 33 misclassifications (12 false positives, 21 false negatives). We sampled 10 of each for manual inspection (Table~\ref{tab:errors}, auto-tagged with heuristic rules and reviewed). Errors concentrate in two categories: \emph{short-answer ambiguity} (very short answers where neither lexical nor NLI evidence is decisive) and \emph{label ambiguity} (answers that are semantically plausible under both labels). This matches the known HaluEval artifact that short answers are frequently hallucinated in the benchmark.

\begin{table}[t]
\centering
\small
\caption{Error taxonomy for the 20 sampled misclassifications (heuristic first-pass tags, reviewed).}
\label{tab:errors}
\begin{tabular}{@{}lcc@{}}
\toprule
\textbf{Category} & \textbf{FP} & \textbf{FN} \\
\midrule
short\_answer\_ambiguity & 6 & 8 \\
label\_ambiguity          & 4 & 2 \\
\bottomrule
\end{tabular}
\end{table}

\subsection{Limitations}

\begin{itemize}
    \item \textbf{Synthetic data.} 85\% of HaluEval is ChatGPT-generated ``sampling-then-filtering'' data; real-world hallucination patterns differ, and our RAGTruth zero-shot result (F1 0.4822) quantifies the gap.
    \item \textbf{Benchmark artifacts.} Length and overlap patterns are highly informative on HaluEval but may be dataset-specific (see ablation and SHAP).
    \item \textbf{Binary labels.} The benchmark reduces hallucination to yes/no; graded severity and partial hallucinations are out of scope.
    \item \textbf{English only.} All features (NER, NLI, embeddings) are English-centric.
    \item \textbf{Single dataset.} The model is trained only on HaluEval QA; domain adaptation is future work.
    \item \textbf{Platt vs.\ isotonic.} Platt did not improve calibration here; isotonic did. The deployed artifact ships Platt (blueprint default).
    \item \textbf{Style-sensitive calibrated score.} In conversational use, full-sentence answers can saturate the calibrated score even when per-claim evidence verdicts say supported; the interface therefore leads with the evidence verdicts and labels the legacy score as secondary.
\end{itemize}
